
\section{\BookRecordIndex{T0200}局部哈密顿纤维化}\label{sec:8-2}

本节把上面引入的概念推广到非平凡纤维丛
\[
  \pi:\widetilde M\longrightarrow\Sigma
\]
的情形。先讨论适当的\BookRecordIndex{T0156}哈密顿纤维化、联络与\BookRecordIndex{T0088}殆复结构概念。

\begin{definition}\label{def:8-2-1}
设 $(M,\omega)$ 是维数为 $2n$ 的紧\BookRecordIndex{T0358}辛流形，
$(\Sigma,\mathrm{dvol}_\Sigma,j_\Sigma)$ 是紧连通黎曼曲面。以 $(M,\omega)$ 为纤维的
$\Sigma$ 上的\emph{\BookRecordIndex{T0200}局部哈密顿纤维化}，是满足下列性质的三元组
$(\widetilde M,\pi,\widetilde\omega)$：
\begin{enumerate}[label=\textup{(\roman*)},leftmargin=3em]
\item $\widetilde M$ 是紧光滑流形，$\pi:\widetilde M\to\Sigma$ 是浸没；
\item $\widetilde\omega$ 是 $\widetilde M$ 上的闭二形式，而且对每个 $z\in\Sigma$，它在
纤维 $M_z:=\pi^{-1}(z)$ 上的限制 $\omega_z$ 都非退化；
\item 对每个 $z\in\Sigma$，纤维 $(M_z,\omega_z)$ 都辛同胚于 $(M,\omega)$；
\item $\pi_*\widetilde\omega^{n+1}=0$，其中
$\pi_*:\Omega^{2n+2}(\widetilde M)\to\Omega^2(\Sigma)$ 表示\BookRecordIndex{T0374}沿纤维积分。
\end{enumerate}
\end{definition}

条件 (i) 断言投影 $\pi:\widetilde M\to\Sigma$ 是局部平凡\BookRecordIndex{T0279}纤维化。定义~\ref{def:8-2-1}
中的二形式 $\widetilde\omega$，在 Guillemin--Lerman--Sternberg \srcite{163} 和
McDuff--Salamon \srcite{277} 第 6 章的术语中是一个\BookRecordIndex{T0224}联络形式。$T\widetilde M$ 中相应的
\BookRecordIndex{T0314}水平分布为
\begin{equation}\label{eq:8-2-1}
  \operatorname{Hor}_{\widetilde x}
  :=\bigl\{\widetilde v\in T_{\widetilde x}\widetilde M
  \mid \widetilde\omega(\widetilde v,\widetilde w)=0,
  \ \forall\widetilde w\in\ker d\pi(\widetilde x)\bigr\}.
\end{equation}
如命题~\ref{prop:8-2-2} 将说明的，$\widetilde\omega$ 为闭这一事实蕴含：绕
$\Sigma$ 中每条可缩环路的\BookRecordIndex{T0164}和乐都是纤维的哈密顿辛同胚。因此，当
$\Sigma=\mathbb{CP}^1$ 时，\BookRecordIndex{T0200}局部哈密顿纤维化在 Lalonde--McDuff \srcite{227} 的意义下
是哈密顿的，也就是说，绕底空间中\emph{每条}环路的\BookRecordIndex{T0164}和乐都是哈密顿辛同胚。不过，
当 $\pi_1(\Sigma)\ne0$ 时未必如此。例如，若
$f:\Sigma\to\Sigma$ 是不与恒等映射同痕的保面积微分同胚，而
\[
  Y_f:=\mathbb R\times\Sigma/(t+1,z)\sim(t,f(z))
\]
是它的\BookRecordIndex{T0381}映射圆，就很容易在四维流形
\[
  X_f:=S^1\times Y_f\longrightarrow\mathbb T^2
\]
上赋予\BookRecordIndex{T0200}局部哈密顿纤维化结构。但它的结构群不能约化到哈密顿辛同胚群
$\operatorname{Ham}(\Sigma)$，因为后者是连通的。条件 (iv) 是归一化条件，它蕴含二形式
$\widetilde\omega$ 由其所定义的联络唯一确定；这就是
Guillemin--Lerman--Sternberg \srcite{163} 的\emph{\BookRecordIndex{T0260}耦合形式}。

\begin{proposition}\label{prop:8-2-2}
设 $(\widetilde M,\pi,\widetilde\omega)$ 是紧黎曼曲面
$(\Sigma,\mathrm{dvol}_\Sigma,j_\Sigma)$ 上以 $(M,\omega)$ 为纤维的\BookRecordIndex{T0200}局部哈密顿纤维化。
则存在 $\Sigma$ 的开覆盖 $\{U_\alpha\}_\alpha$ 与一族局部平凡化
\[
  \widetilde\psi_\alpha:\pi^{-1}(U_\alpha)\longrightarrow U_\alpha\times M,
\]
满足下列条件。
\begin{enumerate}[label=\textup{(\roman*)},leftmargin=3em]
\item 转移映射具有形式
\[
  \widetilde\psi_\beta\circ\widetilde\psi_\alpha^{-1}(z,x)
  =(z,\psi_{\beta\alpha}(z)(x)),
\]
其中 $z\in U_\alpha\cap U_\beta$、$x\in M$，并且
$\psi_{\beta\alpha}:U_\alpha\cap U_\beta\to\operatorname{Diff}(M,\omega)$ 给每个
$z\in U_\alpha\cap U_\beta$ 指派 $(M,\omega)$ 的一个辛同胚
$\psi_{\beta\alpha}(z)$。
\item 对每个 $\alpha$，由
\[
  \widetilde\psi_\alpha^*\widetilde\omega_\alpha=\widetilde\omega
\]
定义的二形式 $\widetilde\omega_\alpha\in\Omega^2(U_\alpha\times M)$ 具有形式
\begin{equation}\label{eq:8-2-2}
  \widetilde\omega_\alpha
  =\omega-d'F_\alpha\wedge ds-d'G_\alpha\wedge dt
   +(\partial_tF_\alpha-\partial_sG_\alpha)\,ds\wedge dt,
\end{equation}
其中 $s+it:U_\alpha\to\mathbb C$ 是 $\Sigma$ 上的全纯坐标卡，$F_\alpha$ 与
$G_\alpha$ 是 $U_\alpha\times M$ 上的实值函数，并且对每个 $z\in U_\alpha$，它们在
$\{z\}\times M$ 上的平均值为零；$d'$ 表示对第二个因子取微分。
\end{enumerate}
\end{proposition}

\begin{proof}
因为 $\widetilde\omega$ 是闭的，由 Cartan 公式，沿光滑曲线
$\gamma:[0,1]\to\Sigma$ 的\BookRecordIndex{T0263}平行移动保持纤维上的\BookRecordIndex{T0363}辛形式。更明确地，若
\[
  u_t:M_{\gamma(0)}\longrightarrow M_{\gamma(t)}\subset\widetilde M
\]
是沿 $\gamma$ 的水平\BookRecordIndex{T0263}平行移动族，\sourcefootnote{原书把 $u_t$ 的定义域印作
$\Sigma$；但此处 $u_t^*\widetilde\omega$、$du_t$ 以及沿 $\gamma$ 的纤维间\BookRecordIndex{T0263}平行移动都要求
定义域为起始纤维 $M_{\gamma(0)}$，故据此补正。}则
\[
  \partial_tu_t^*\widetilde\omega=d\alpha_t,
  \qquad
  \alpha_t:=\widetilde\omega(\partial_tu_t,du_t\,\mathord\cdot)=0.
\]
（另见引理~\ref{lem:8-1-1} 前的讨论。）因此，沿适当的路径族作\BookRecordIndex{T0263}平行移动，就可以构造
满足 (i) 中条件的平凡化。

在可缩开集 $U_\alpha\subset\Sigma$ 上给定这样的平凡化，并在其上取\BookRecordIndex{T0149}共形坐标 $s+it$。
考察底空间中的环路可知，向量场 $\partial_s$、$\partial_t$ 到 $U_\alpha\times M$ 的\BookRecordIndex{T0317}水平
提升 $(\partial_s,-X_{s,t})$、$(\partial_t,-Y_{s,t})$ 沿 $M$ 是哈密顿的。这就是说，
$X_{s,t}$ 与 $Y_{s,t}$ 对所有 $s,t$ 都是 $M$ 上的\BookRecordIndex{T0160}哈密顿向量场，因而分别由 $M$ 上的函数族
$F_\alpha(s,t,\mathord\cdot)$ 与 $G_\alpha(s,t,\mathord\cdot)$ 生成。（详细证明见
\srcite{277}，定理 6.21。）这些函数可以选成在 $M$ 上平均值为零。由\BookRecordIndex{T0314}水平分布的定义，
\[
  \widetilde\omega_\alpha
  =\omega-d'F_\alpha\wedge ds-d'G_\alpha\wedge dt+f\,ds\wedge dt.
\]
由于 $\widetilde\omega_\alpha$ 是闭的，得到
$d'f=\partial_td'F_\alpha-\partial_sd'G_\alpha$；这把 $f$ 确定到相差一个 $s+it$ 的函数。
最后可以验证，定义~\ref{def:8-2-1} 中的耦合条件 (iv)，等价于 $f$ 在每条纤维
$\{s+it\}\times M$ 上平均值为零。因此
$f=\partial_tF_\alpha-\partial_sG_\alpha$，正如式~\eqref{eq:8-2-2} 所示。
\end{proof}

\paragraph{曲率。}
$\widetilde\omega$ 所确定的 $\widetilde M$ 上联络之曲率，是 $\Sigma$ 上取值于某个丛的
二形式 $R\,\mathrm{dvol}_\Sigma$；这个丛在 $z$ 处的纤维为
$C_0^\infty(M_z)$，即 $M_z$ 上关于体积形式 $\omega_z^n$ 平均值为零的光滑函数空间。
这与通常的曲率概念一致，因为 $C_0^\infty(M)$（等价地，$M$ 上\BookRecordIndex{T0160}哈密顿向量场所成的空间）
是纤维丛结构群的李代数。我们仍把曲率形式写成 $R\,\mathrm{dvol}_\Sigma$，于是
$R:\widetilde M\to\mathbb R$ 是光滑函数。

曲率在任意一对向量场 $v,w\in\operatorname{Vect}(\Sigma)$ 上的值，是它们的\BookRecordIndex{T0317}水平提升
$v^\#,w^\#\in\operatorname{Vect}(\widetilde M)$ 的李括号 $[v^\#,w^\#]$ 的竖直部分
（见 \srcite{277}，注 6.31）；这里竖直部分定义为 $[v^\#,w^\#]$ 在沿\BookRecordIndex{T0314}水平分布的投影
\[
  \Pi:T\widetilde M\longrightarrow T^{\mathrm{Vert}}\widetilde M
\]
下的像。因此，对每个 $z\in\Sigma$，$[v^\#,w^\#]^{\mathrm{vert}}$ 在 $M_z$ 上的限制，
是函数
\[
  R(z)\,\mathrm{dvol}_\Sigma(v(z),w(z))\in C_0^\infty(M_z)
\]
的\BookRecordIndex{T0160}哈密顿向量场。取 $v:=\partial_s$、$w:=\partial_t$，可以证明在
命题~\ref{prop:8-2-2} 的局部平凡化中，曲率形式由
\[
  R_\alpha\lambda_\alpha
  =\partial_sG_\alpha-\partial_tF_\alpha+\{F_\alpha,G_\alpha\}
\]
给出。这里 $R_\alpha:U_\alpha\times M\to\mathbb R$ 由
$R_\alpha\circ\widetilde\psi_\alpha=R$ 定义，而
$\lambda_\alpha:U_\alpha\to\mathbb R$ 通过
$\mathrm{dvol}_\Sigma=\lambda_\alpha\,ds\wedge dt$ 确定 $\Sigma$ 上的体积形式。因此，
由下面习题~\ref{ex:8-2-3} 中的式~\eqref{eq:8-2-4}，有
\begin{equation}\label{eq:8-2-3}
  R\,\mathrm{dvol}(v,w)=-\widetilde\omega(v^\#,w^\#)
\end{equation}
其中 $v,w$ 是 $\Sigma$ 上任意一对向量。（更多细节见 \srcite{277}，第 6 章。\footnote{尽管
这里的整体符号约定与 \srcite{277} 相同，本节称为 $F,G$ 的函数，分别相当于
\srcite{277} 习题 6.30 中的 $-H_i,-H_j$。此外，\srcite{277} 引理 6.28 有一个符号错误；
$d\tau$ 的公式中所有项都应取正号。}）

\begin{exercise}\label{ex:8-2-3}
沿用命题~\ref{prop:8-2-2} 的记号。证明，$\widetilde\omega_\alpha$ 在水平向量场
\[
  \partial_s^\#=(\partial_s,-X_{F_\alpha}),
  \qquad
  \partial_t^\#=(\partial_t,-X_{G_\alpha})
\]
上的值为
\begin{equation}\label{eq:8-2-4}
  \widetilde\omega_\alpha(\partial_s^\#,\partial_t^\#)
  =-\partial_sG_\alpha+\partial_tF_\alpha-\{F_\alpha,G_\alpha\},
\end{equation}
其中 $\{\mathord\cdot,\mathord\cdot\}$ 表示 $M$ 上的 Poisson 括号。
\end{exercise}

\begin{exercise}\label{ex:8-2-4}
设 $\psi_{\beta\alpha}:U_\alpha\cap U_\beta\to\operatorname{Diff}(M,\omega)$ 如
命题~\ref{prop:8-2-2} 所示。证明，这些辛同胚族由\BookRecordIndex{T0160}哈密顿向量场生成。（辛同胚
$\psi_{\beta\alpha}(z)$ 本身未必是哈密顿的，但它们彼此\BookRecordIndex{T0158}哈密顿同痕。）在
$U_\alpha\cap U_\beta$ 上给定\BookRecordIndex{T0149}共形坐标 $s+it$，用条件
\[
  \partial_s\psi_{\beta\alpha}
  +X_{F_{\beta\alpha}}\circ\psi_{\beta\alpha}=0,
  \qquad
  \partial_t\psi_{\beta\alpha}
  +X_{G_{\beta\alpha}}\circ\psi_{\beta\alpha}=0
\]
以及如下归一化条件定义函数
\[
  F_{\beta\alpha},G_{\beta\alpha}:
  (U_\alpha\cap U_\beta)\times M\longrightarrow\mathbb R:
\]
对每个 $z\in U_\alpha\cap U_\beta$，$F_{\beta\alpha}$ 与 $G_{\beta\alpha}$ 在
$\{z\}\times M$ 上平均值为零。证明在 $U_\alpha\cap U_\beta$ 上
\[
  F_\alpha=(F_\beta-F_{\beta\alpha})\circ\psi_{\beta\alpha},
  \qquad
  G_\alpha=(G_\beta-G_{\beta\alpha})\circ\psi_{\beta\alpha}.
\]
由此\BookRecordIndex{T0326}推出
\[
  \partial_sG_\alpha-\partial_tF_\alpha+\{F_\alpha,G_\alpha\}
  =\bigl(\partial_sG_\beta-\partial_tF_\beta+\{F_\beta,G_\beta\}\bigr)
   \circ\psi_{\beta\alpha}.
\]
\end{exercise}

\begin{exercise}\label{ex:8-2-5}
设 $\kappa:\Sigma\to\mathbb R$ 是光滑函数，并且
$\kappa\circ\pi-R:\widetilde M\to\mathbb R$ 处处不为零。证明二形式
\[
  \widetilde\omega_\kappa
  :=\widetilde\omega+\pi^*(\kappa\,\mathrm{dvol}_\Sigma)
\]
是 $\widetilde M$ 上的\BookRecordIndex{T0363}辛形式。（见习题~\ref{ex:8-1-3}。）
\end{exercise}

\paragraph{\BookRecordIndex{T0088}殆复结构。}
我们关心\BookRecordIndex{T0200}局部哈密顿纤维化的伪全纯截面。为此，需要考虑如下意义下适配于\BookRecordIndex{T0279}纤维化的\BookRecordIndex{T0088}殆复结构。
设 $(\widetilde M,\pi,\widetilde\omega)$ 是紧黎曼曲面
$(\Sigma,\mathrm{dvol}_\Sigma,j_\Sigma)$ 上的\BookRecordIndex{T0200}局部哈密顿纤维化，其纤维辛同胚于
$(M,\omega)$。

\begin{definition}\label{def:8-2-6}
若 $(\widetilde M,\widetilde\omega)$ 上的\BookRecordIndex{T0088}殆复结构 $\widetilde J$ 满足下列条件，就称它
\emph{与\BookRecordIndex{T0279}纤维化\BookRecordIndex{T0348}相容}。
\begin{enumerate}[label=\textup{(\roman*)},leftmargin=3em]
\item 投影 $\pi:\widetilde M\to\Sigma$ 是全纯的，即
\[
  d\pi\circ\widetilde J=j_\Sigma\circ d\pi.
\]
\item 对每个 $z\in\Sigma$，$\widetilde J$ 在
$M_z=\pi^{-1}(z)$ 上的限制 $J_z$ 都被 $\omega_z$ \BookRecordIndex{T0372}驯服。
\item \BookRecordIndex{T0314}水平分布 $\operatorname{Hor}\subset T\widetilde M$ 在 $\widetilde J$ 下不变。
\end{enumerate}
$\widetilde M$ 上与\BookRecordIndex{T0279}纤维化\BookRecordIndex{T0348}相容的\BookRecordIndex{T0088}殆复结构所成的集合记作
$\mathcal J(\widetilde M,\pi,\widetilde\omega)$。
\end{definition}

有时我们只关心 $\widetilde J$ 的竖直部分。因此作如下定义。注意我们的约定：
$\widetilde J$ 表示整个 $T\widetilde M$ 上的\BookRecordIndex{T0088}殆复结构，而
$J=\{J_z\}$ 只涉及它的竖直子丛。

\begin{definition}\label{def:8-2-7}
$\widetilde M$ 上一个\emph{被 $\widetilde\omega$ \BookRecordIndex{T0372}驯服的\BookRecordIndex{T0306}竖直殆复结构}，是\BookRecordIndex{T0308}竖直切丛
\[
  T^{\mathrm{Vert}}\widetilde M:=\ker d\pi\subset T\widetilde M
\]
上的自同构 $J$，满足 $J^2=-\mathbf1$，而且对每个非零向量
$\widetilde v\in T^{\mathrm{Vert}}\widetilde M$
\sourcefootnote{原书把严格不等式写成对每个竖直向量都成立，因而错误地包含了零向量；
\BookRecordIndex{T0372}驯服条件只对非零向量要求严格正，故补上这一限制。}都有
\[
  \widetilde\omega(\widetilde v,J\widetilde v)>0.
\]
这类\BookRecordIndex{T0088}殆复结构所成的空间记作
$\mathcal J^{\mathrm{Vert}}(\widetilde M,\pi,\widetilde\omega)$。
\end{definition}

等价地，一个被 $\widetilde\omega$ \BookRecordIndex{T0372}驯服的\BookRecordIndex{T0306}竖直殆复结构，可以描述为纤维 $M_z$ 上\BookRecordIndex{T0088}殆复结构
的光滑族 $J=\{J_z\}_{z\in\Sigma}$，并且每个 $z\in\Sigma$ 都满足
$J_z\in\mathcal J_\tau(M_z,\omega_z)$。若
$J\in\mathcal J^{\mathrm{Vert}}(\widetilde M,\pi,\widetilde\omega)$，则式~\eqref{eq:2-1-1}
在 $T^{\mathrm{Vert}}\widetilde M$ 上定义一个内积。

\begin{lemma}\label{lem:8-2-8}
对每个 $J\in\mathcal J^{\mathrm{Vert}}(\widetilde M,\pi,\widetilde\omega)$，存在唯一的
$\widetilde J\in\mathcal J(\widetilde M,\pi,\widetilde\omega)$ 与\BookRecordIndex{T0279}纤维化\BookRecordIndex{T0348}相容，并且在竖直
子丛 $T^{\mathrm{Vert}}\widetilde M$ 上限制为 $J$。
\end{lemma}

\begin{proof}
这直接来自定义。
\end{proof}

上述概念与第~\ref{sec:8-1} 节引入的概念一致。因此，在命题~\ref{prop:8-2-2} 的局部
平凡化中，\BookRecordIndex{T0088}殆复结构 $\widetilde J$ 由公式~\eqref{eq:8-1-6} 给出，而相应的
$\widetilde J$-全纯曲线由式~\eqref{eq:8-1-7} 给出。现在令
$\widetilde J\in\mathcal J(\widetilde M,\pi,\widetilde\omega)$，并令
$\widetilde u:\Sigma\to\widetilde M$ 为 $\widetilde J$-全纯截面。$\widetilde u$ 的
\emph{\BookRecordIndex{T0307}竖直能量}定义为微分
\[
  d\widetilde u:T\Sigma\longrightarrow\widetilde u^*T\widetilde M
\]
的竖直部分
$\Pi d\widetilde u\in\Omega^1(\Sigma,\widetilde u^*T^{\mathrm{Vert}}\widetilde M)$
的 $L^2$-范数平方的一半。这里投影
$\Pi:T\widetilde M\to T^{\mathrm{Vert}}\widetilde M$ 由
\begin{equation}\label{eq:8-2-5}
  \Pi(\widetilde x)\widetilde v
  :=\widetilde v-\bigl(d\pi(\widetilde x)\widetilde v\bigr)^\#
\end{equation}
给出，其中 $\widetilde v\in T_{\widetilde x}\widetilde M$，而
$\zeta^\#\in T_{\widetilde x}\widetilde M$ 表示
$\zeta\in T_{\pi(\widetilde x)}\Sigma$ 的\BookRecordIndex{T0317}水平提升。因此
\[
  E^{\mathrm{Vert}}(\widetilde u)
  :=\frac12\int_\Sigma\lvert\Pi d\widetilde u\rvert_{\widetilde J}^2
  \,\mathrm{dvol}_\Sigma.
\]
\BookRecordIndex{T0307}竖直能量为零，当且仅当 $\widetilde u$ 是 $\widetilde M$ 的\BookRecordIndex{T0315}水平截面。\BookRecordIndex{T0307}竖直能量恒等式涉及
$\widetilde\omega$ 所确定联络的曲率，并且具有与前面相同的形式；参见
式~\eqref{eq:8-1-9}。

\begin{lemma}\label{lem:8-2-9}
设 $(\widetilde M,\pi,\widetilde\omega)$ 是 $\Sigma$ 上的\BookRecordIndex{T0200}局部哈密顿纤维化，
$\widetilde J\in\mathcal J(\widetilde M,\pi,\widetilde\omega)$，且
$\widetilde u:\Sigma\to\widetilde M$ 是 $\widetilde J$-全纯截面。则
\[
  E^{\mathrm{Vert}}(\widetilde u)
  =\int_\Sigma\widetilde u^*\widetilde\omega
   +\int_\Sigma R(\widetilde u)\,\mathrm{dvol}_\Sigma.
\]
\end{lemma}

\begin{proof}
在带坐标 $z=s+it$ 的局部坐标卡 $U_\alpha\subset\Sigma$ 中，局部能量满足下式。
\sourcefootnote{原书把第一行 $E^{\mathrm{Vert}}(\widetilde u;U_\alpha)$ 的积分域误印为
$\Sigma$；该量是 $U_\alpha$ 上的局部能量，且随后各行均在 $U_\alpha$ 上积分，故补正为
$U_\alpha$。}
\begin{align*}
  E^{\mathrm{Vert}}(\widetilde u;U_\alpha)
  &=\frac12\int_{U_\alpha}
    \lvert\Pi d\widetilde u\rvert_{\widetilde J}^2\,
    \mathrm{dvol}_\Sigma\\
  &=\int_{U_\alpha}
    \widetilde\omega(\Pi\partial_s\widetilde u,
                     \Pi\widetilde J\partial_s\widetilde u)
    \,\mathrm{dvol}_\Sigma\\
  &=\int_{U_\alpha}
    \left(
      \widetilde\omega(\Pi\partial_s\widetilde u+\partial_s^\#,
                       \Pi\partial_t\widetilde u+\partial_t^\#)
      -\widetilde\omega(\partial_s^\#,\partial_t^\#)
    \right)\mathrm{dvol}_\Sigma\\
  &=\int_{U_\alpha}
    \left(\widetilde\omega(\partial_s\widetilde u,
                            \partial_t\widetilde u)
          +R(\widetilde u)\right)\mathrm{dvol}_\Sigma.
\end{align*}
第一行使用了 $\Pi$ 与 $\widetilde J$ 可交换这一事实；第二行使用了按定义向量的水平部分
与竖直部分关于 $\widetilde\omega$ 正交；最后一行使用式~\eqref{eq:8-2-3}。也可以利用
命题~\ref{prop:8-2-2} 的局部平凡化，从式~\eqref{eq:8-1-9} \BookRecordIndex{T0326}推出这一公式。
\end{proof}

\begin{remark}[\BookRecordIndex{T0224}联络形式族]\label{rem:8-2-10}
给定\BookRecordIndex{T0224}联络形式 $\widetilde\omega$ 与\BookRecordIndex{T0306}竖直殆复结构 $J=\{J_z\}$，由
引理~\ref{lem:8-2-8}，存在唯一的、与 $\widetilde\omega$ \BookRecordIndex{T0348}相容的\BookRecordIndex{T0088}殆复结构
$\widetilde J$，从而有唯一确定的\BookRecordIndex{T0209}柯西--黎曼方程。为了证明\BookRecordIndex{T0408}正则性结果，需要考虑一族方程。
得到这种方程族的一种方式是扰动 $J$。不过，另一种很有用的办法是固定 $J$，转而扰动
\BookRecordIndex{T0224}联络形式 $\widetilde\omega$。

因此，我们要考虑 $\widetilde M\to\Sigma$ 上的\BookRecordIndex{T0224}联络形式族，它们在各纤维上都等于给定形式
$\widetilde\omega$。每个这样的\BookRecordIndex{T0224}联络形式都可以写成 $\widetilde\omega$ 的\BookRecordIndex{T0402}正合扰动
\[
  \widetilde\omega_H:=\widetilde\omega-d\widetilde\sigma_H,
\]
其中一形式 $\widetilde\sigma_H$ 在纤维上消失。此外，由命题~\ref{prop:8-2-2}(ii)，
$\widetilde\omega$ 和 $\widetilde\omega_H$ 在局部坐标中都由形如~\eqref{eq:8-2-2} 的表达式
给出。因此，实值一形式 $\widetilde\sigma_H$ 必须具有第~\ref{sec:8-1} 节所考察的类型。
换言之，它来自 $\Sigma$ 上的一个一形式 $H$，这个一形式取值于如下函数丛：其在
$z\in\Sigma$ 处的纤维是 $C_0^\infty(M_z)$。若后一形式记作
\[
  T_z\Sigma\longrightarrow C_0^\infty(M_z):\quad
  \zeta\longmapsto H_\zeta,
\]
则 $\widetilde M$ 上相应的一形式 $\widetilde\sigma_H$ 由
\[
  \widetilde\sigma_H(\widetilde x;\widetilde v)=H_\zeta(\widetilde x),
  \qquad
  \zeta:=d\pi_{\widetilde x}(\widetilde v),
\]
给出，其中 $\widetilde v\in T_{\widetilde x}\widetilde M$。在此记号下，与\BookRecordIndex{T0224}联络形式
$\widetilde\omega-d\widetilde\sigma_H$ 对应的 $\widetilde x\in\widetilde M$ 处水平子空间为
\[
  \operatorname{Hor}_{H;\widetilde x}
  =\bigl\{\widetilde v-X_{H_\zeta}(\widetilde x)
  \mid \widetilde v\in\operatorname{Hor}_{\widetilde x},\ 
  \zeta=d\pi_{\widetilde x}(\widetilde v)\bigr\},
\]
其中 $\operatorname{Hor}_{\widetilde x}\subset T_{\widetilde x}\widetilde M$ 是与
$\widetilde\omega$ 对应的水平子空间（由式~\eqref{eq:8-2-1} 给出），而对
$\zeta\in T_z\Sigma$，$X_{H_\zeta}\in\operatorname{Vect}(M_z,\omega_z)$ 表示哈密顿函数
$H_\zeta:M_z\to\mathbb R$ 所对应的\BookRecordIndex{T0160}哈密顿向量场。由此得到
\[
  \widetilde J_H\widetilde v
  =\widetilde J\widetilde v
   +JX_{H_\zeta}(\widetilde x)-X_{H_{j\zeta}}(\widetilde x),
  \qquad
  \zeta:=d\pi_{\widetilde x}(\widetilde v),
\]
其中 $\widetilde v\in T_{\widetilde x}\widetilde M$，$j:=j_\Sigma$。
\end{remark}

\paragraph{哈密顿辛同胚环路。}
本节最后在（局部平凡的）\BookRecordIndex{T0156}哈密顿纤维化
\[
  \pi:\widetilde M\longrightarrow S^2
\]
上显式构造\BookRecordIndex{T0224}联络形式。最重要的一点是注~\ref{rem:8-2-11}(iv)，第~\ref{sec:9-6} 节将会用到它。
由命题~\ref{prop:8-2-2}，纤维丛 $\pi:\widetilde M\to S^2$ 在每个半球上都可平凡化；
因此，可以把总空间 $\widetilde M$ 表为圆盘上两个平凡丛之并，它们沿公共边界\BookRecordIndex{T0254}粘合。
这个\BookRecordIndex{T0257}粘合映射是哈密顿辛同胚的一条环路。反过来，给定这样一条环路，可以如下构造
$S^2$ 上的\BookRecordIndex{T0279}纤维化。

令 $(M,\omega)$ 为紧\BookRecordIndex{T0358}辛流形，并令
\[
  \mathbb R/\mathbb Z\longrightarrow\operatorname{Diff}(M,\omega):
  \quad t\longmapsto\psi_t
\]
为哈密顿辛同胚的一条环路。与环路 $\psi=\{\psi_t\}$ 对应的\BookRecordIndex{T0156}哈密顿纤维化
\[
  M\lhook\joinrel\longrightarrow\widetilde M_\psi
  \xrightarrow{\ \pi\ }S^2
\]
由环路 $\psi$ \BookRecordIndex{T0254}粘合两份 $\mathbb C\times M$ 定义：
\begin{equation}\label{eq:8-2-6}
  \widetilde M_\psi:=\{\pm1\}\times\mathbb C\times M/\sim,
\end{equation}
其中
\[
  (1,e^{2\pi(s+it)},x)
  \sim(-1,e^{2\pi(-s-it)},\psi_t^{-1}(x)).
\]
投影 $\pi:\widetilde M_\psi\to S^2=\mathbb C\cup\{\infty\}$ 由
$\pi([1,z,x]):=z$ 给出。我们常在 $\mathbb C\setminus\{0\}$ 上使用极坐标 $s,t$，令
$z=e^{2\pi(s+it)}$。在此记号下，$\widetilde M_\psi$ 的一个截面是一对
$\widetilde u=(u^+,u^-)$，其中
$u^\pm:\mathbb R\times\mathbb R/\mathbb Z\to M$ 是光滑函数，满足
\[
  u^+(s,t)=\psi_t\bigl(u^-(-s,-t)\bigr),
\]
并且映射
\[
  \mathbb C\setminus\{0\}\longrightarrow M:\quad
  e^{2\pi(s+it)}\longmapsto u^\pm(s,t)
\]
都能光滑延拓过原点。

\begin{remark}\label{rem:8-2-11}
\begin{enumerate}[label=\textup{(\roman*)},leftmargin=3em]
\item \BookRecordIndex{T0224}联络形式 $\widetilde\omega\in\Omega^2(\widetilde M_\psi)$ 可以表示为一对二形式
$\omega^\pm\in\Omega^2(\mathbb C\times M)$。在极坐标中把它们写成
\[
  \omega^\pm
  =\omega-d'F^\pm\wedge ds-d'G^\pm\wedge dt
   +(\partial_tF^\pm-\partial_sG^\pm)\,ds\wedge dt,
\]
其中 $d'$ 表示 $M$ 上的微分，而且函数
\[
  e^{2\pi(s+it)}\longmapsto
  e^{2\pi(-s-it)}\bigl(F^\pm_{s,t}+iG^\pm_{s,t}\bigr)
\]
都能光滑延拓过原点。函数 $F^\pm$ 与 $G^\pm$ 满足关系
\[
  F^+_{s,t}\circ\psi_t+F^-_{-s,-t}=0,
  \qquad
  G^+_{s,t}\circ\psi_t+G^-_{-s,-t}=-H_t\circ\psi_t,
\]
其中哈密顿函数 $H_t=H_{t+1}:M\to\mathbb R$ 通过
$\partial_t\psi_t=X_{H_t}\circ\psi_t$ 生成哈密顿环路 $\psi$，并且对每个 $t$ 都满足
\[
  \int_MH_t\omega^n=0.
\]

\item 令 $\sigma\in\Omega^2(S^2)$ 为 $S^2=\mathbb C\cup\{\infty\}$ 上总体积为
$\pi$ 的标准体积形式。在极坐标
\[
  \mathbb R\times\mathbb R/\mathbb Z\longrightarrow\mathbb C:\quad
  (s,t)\longmapsto e^{2\pi(s+it)}
\]
下，这个体积形式为
\[
  \sigma=\frac{4\pi^2e^{4\pi s}}{(1+e^{4\pi s})^2}\,ds\wedge dt.
\]

\item 如 (i) 中的\BookRecordIndex{T0224}联络形式
$\widetilde\omega\in\Omega^2(\widetilde M_\psi)$ 的曲率，是 $S^2$ 上取值于纤维上平均值
为零的函数丛的二形式 $R_{\widetilde\omega}\sigma$。这里函数
$R=R_{\widetilde\omega}:\widetilde M_\psi\to\mathbb R$ 在两个平凡化中由
\[
  R^\pm(s,t,x)
  :=\frac{(1+e^{4\pi s})^2}{4\pi^2e^{4\pi s}}
  \left(
    \partial_sG^\pm_{s,t}(x)-\partial_tF^\pm_{s,t}(x)
    +\{F^\pm_{s,t},G^\pm_{s,t}\}(x)
  \right)
\]
给出。

\item 令 $\kappa\in\Omega^0(S^2)$，并记
\[
  \kappa^\pm(s,t):=\kappa(e^{\pm2\pi(s+it)}),
  \qquad
  \omega^\pm_\kappa:=\omega^\pm+\kappa^\pm\pi^*\sigma.
\]
公式
\[
  \frac{(\omega^\pm_\kappa)^{n+1}}{(n+1)!}
  =(\kappa^\pm-R^\pm_{\widetilde\omega})\,
    \sigma\wedge\frac{\omega^n}{n!}
\]
表明，可以由二形式
$\widetilde\omega_\kappa:=\widetilde\omega+\pi^*(\kappa\sigma)$，通过\BookRecordIndex{T0374}沿纤维积分恢复函数
$\kappa$：
\[
  \kappa\sigma
  :=\frac1{\operatorname{Vol}(M)}
    \int_{\mathrm{fiber}}
    \frac{\widetilde\omega_\kappa^{n+1}}{(n+1)!}.
\]
这个公式还表明，$\widetilde\omega_\kappa$ 是\BookRecordIndex{T0363}辛形式，当且仅当函数
$\kappa\circ\pi-R_{\widetilde\omega}$ 处处不为零。约定
\[
  \widetilde\omega_\kappa^{n+1}>0
  \quad\Longleftrightarrow\quad
  \kappa\circ\pi>R_{\widetilde\omega},
  \qquad
  \widetilde\omega_\kappa^{n+1}<0
  \quad\Longleftrightarrow\quad
  \kappa\circ\pi<R_{\widetilde\omega}.
\]
\end{enumerate}
\end{remark}
